% Standalone input section; requires amsmath, amssymb, amsthm in parent document.
\Needspace{10\baselineskip}
\section{A finite certificate for the conjugate-dual congruence term}
\label{sec:congruence-smith}

The following is elementary linear algebra over a discrete valuation ring. Its
value here is diagnostic: it identifies exactly when a hypothesis of the
upstream twisting theorem can hold, and provides a finite certificate when it
does. The structural incoming report already uses this viewpoint
\cite{IncomingStructural}; we make no claim of historical novelty for the
linear algebra.

Let $O$ be a discrete valuation ring with uniformizer $\pi$ and fraction field
$K$. Let $T,U$ be finite free $O$-modules with actions of a group $G$. Put
$V=\operatorname{Hom}_O(T,U)$, of rank $m$, and let $R\subseteq V^*$ be the
$O$-submodule generated by all coordinate equations of
\[
 \rho_U(g)f-f\rho_T(g)=0\qquad(g\in G).
\]
Write
\[
 Q:=V^*/R\simeq O^d\oplus\bigoplus_i O/(\pi^{a_i}),\qquad a_i>0.
\]
Then for every $k\geq1$ there is an $O$-linear identification
\[
 \operatorname{Hom}_G(T/\pi^kT,U/\pi^kU)
 \simeq\operatorname{Hom}_O(Q,O/\pi^k),
\]
and consequently the exact formula
\begin{equation}
 \operatorname{length}_O\operatorname{Hom}_G(T/\pi^kT,U/\pi^kU)
   =dk+\sum_i\min(k,a_i).
 \label{eq:congruence-smith-length}
\end{equation}

Indeed, a vector of $V/\pi^kV$ is a functional on $V^*$ with values in
$O/\pi^k$; equivariance says exactly that it annihilates $R$. A free summand
of $Q$ contributes length $k$, whereas $O/\pi^a$ contributes the
$\pi^a$-torsion in $O/\pi^k$, of length $\min(a,k)$. Furthermore,
\[
 d=\dim_K\operatorname{Hom}_G(T\otimes_OK,U\otimes_OK).
\]
It follows that the supremum of the left-hand side of
\eqref{eq:congruence-smith-length} is finite if and only if this rational Hom
space is zero. In that case the supremum is $\sum_i a_i$, reached for
$k\geq\max_i a_i$ when torsion factors are present, and immediately if
there are none. Nonzero ideals of a DVR are powers of its maximal ideal,
so the same conclusion holds for the supremum over arbitrary $a\in O\setminus\{0\}$.

\paragraph{Finite certificates and their limitation.}
Noetherianity implies that finitely many of the original coordinate equations
generate $R$. Finding such a complete set effectively is a separate issue.
There is, however, a useful sufficient certificate: choose any finite relation
matrix $M$ of full column rank $m$. If its Smith exponents are
$b_1,\ldots,b_m\geq0$, the true equivariant Hom is a submodule of the kernel
of $M$ modulo $\pi^k$, whence
\[
 \operatorname{length}_O\operatorname{Hom}_G(T/\pi^kT,U/\pi^kU)
 \leq\sum_{i=1}^m\min(k,b_i)\leq\sum_{i=1}^m b_i.
\]
Any nonzero maximal minor $\Delta$ gives the coarser bound
$\sum b_i\leq v_\pi(\Delta)$. Thus a finite sample can certify an upper
bound without a claim that it generates all Galois relations. It cannot by
itself certify that a nonzero kernel persists after the omitted relations.

\paragraph{Exact correspondence with \texttt{grh-lean}.}
Here take $O$ to be the integer ring of a finite extension of
$\mathbb Q_\ell$, as in the source interface.
In \texttt{New.Selmer.conjDualHom}, take $G=\operatorname{Gal}(\Omega/E)$,
$T=O^n$, $A_g=\rho(g)$, $c(g)=\texttt{conjGal}(g)$, and let $\chi$ denote
the $\ell$-adic cyclotomic character. The source equation for $\eta$ is
\[
 \eta(A_g u)(v)=\eta(u)\bigl(\chi(c(g))A_{c(g)^{-1}}v\bigr).
\]
The source's \texttt{CharacterModule} is the additive character dual with
values in $\mathbb Q/\mathbb Z$. Since the trace-dual lattice
$O^\vee=\operatorname{Hom}_{\mathbb Z_\ell}(O,\mathbb Z_\ell)$ is an
invertible, hence free rank-one, $O$-module, choosing a generator identifies
$a^{-1}O^\vee/O^\vee$ with $O/a$. Thus $\eta$ is represented, noncanonically,
by a balanced bilinear pairing matrix $B$ modulo $a$. Its relations are
exactly
\begin{equation}
 A_g^{\mathsf T}B A_{c(g)}-\chi(g)B=0\pmod a,
 \label{eq:conj-dual-relations}
\end{equation}
using $\chi(c(g))=\chi(g)$. The relation module therefore has $m=n^2$.
Its rational kernel is
$\operatorname{Hom}_G(T_K,T_K^{\vee,c}(1))$.
Accordingly, the upstream \texttt{congruenceTerm} is finite precisely when
this rational Hom vanishes, and then its value is the sum of the Smith
exponents of all the relations. This concerns only the representation term;
it does not identify the upstream cocycle model with a standard Selmer group.

\paragraph{Ramification.}
Let $O'/O$ have relative ramification index $e$ and transport the \emph{same}
representations by scalar extension. Flatness gives $Q'=Q\otimes_OO'$;
the free rank remains $d$ and the elementary exponents become $ea_i$.
Hence the length at $\pi'^k$ is $dk+\sum_i\min(k,ea_i)$.
If $e_0=v_\pi(\ell)$, the normalization in the source divides $O$-length by
$e_0$ and $O'$-length by $ee_0$. The finite normalized supremum is invariant:
\[
 \frac{e\sum_i a_i}{ee_0}=\frac{\sum_i a_i}{e_0}.
\]
At corresponding moduli $\pi^k$ and $\pi'^{ek}$, normalized lengths agree
term by term. Coefficient extension alone cannot remove a rational
conjugate-dual obstruction.

This does not license a vanishing argument for \emph{varying} twists and
coefficient fields. The nonzero module $O'/\pi'$ has normalized length
$1/(ee_0)$, tending to zero as $e$ grows. A fixed-ring bound strictly below
$1/e_0$ forces vanishing, but a bound with the ring and twist depending on
the error tolerance does not. This quantifier distinction is required when
using \texttt{exists\_anticyclotomic\_twist}.

\paragraph{Exact examples and reproducibility.}
The accompanying checker uses only Python's standard library. It obtains
integer Smith factors from gcds of all minors and independently enumerates
small modular kernels. For a rank-one relation $5^a f=0$ over $\mathbb Z_5$,
the length is $\min(k,a)$; for the zero relation it is $k$. Thus residual
isomorphism alone is not the obstruction: rational equivariant Hom is.
For the abstract cyclic example $A=C=\operatorname{diag}(1,26)$,
$\chi=126$, relation~\eqref{eq:conj-dual-relations} has $5$-adic Smith
exponents $(2,2,2,3)$ and supremum $9$. Restriction to symmetric and
alternating matrices gives exponents $(2,2,3)$ and $(2)$, respectively.
These are illustrative matrix certificates, not arithmetic Galois-image
certificates. In residue characteristic two, an alternating constraint must
include zero diagonal entries in addition to off-diagonal skew symmetry.
